% concatenated from ReciprocalDivisorV187.tex
\documentclass[12pt]{article}

\usepackage{amsmath,amssymb,amsfonts}
\usepackage{amsthm}
\usepackage{enumerate}
\usepackage{mathrsfs}

\allowdisplaybreaks%

{\theoremstyle{plain}%
 \newtheorem{theorem}{Theorem}
 \newtheorem{corollary}{Corollary}
 \newtheorem{lemma}{Lemma}% 
}
{\theoremstyle{remark}
\newtheorem{remark}{Remark}
}
{\theoremstyle{definition}
\newtheorem{definition}{Definition}
\newtheorem{example}{Example}
}

\begin{document}

\begin{center}
 {\Large On divisor sums due to Erd\H os and Ramanujan}

 \ 

{\sc John M. Campbell}

\vspace{0.1in}

{\footnotesize Department of Mathematics and Statistics}

{\footnotesize Dalhousie University}

{\footnotesize Halifax, NS B3H 4R2}

{\footnotesize Canada}

\vspace{0.1in}

{\footnotesize {\tt jh241966@dal.ca}}

\vspace{0.1in}

\end{center}

\begin{abstract}
 Let $d(n)$ denote the number of divisors of a positive integer $n$. A classical problem in analytic number theory is given by the 
 asymptotic behavior of the divisor sum $\sum_{n \leq x} \frac{1}{d(n)}$, with Ramanujan having introduced an asymptotic formula for 
 this sum with an explicit evaluation for the constant $A_1$ for the leading term $A_1 \frac{x}{\sqrt{\log x}}$. Gabdullin et al.\ recently 
 considered a hybrid of this problem and the Titchmarsh divisor problem concerning $\sum_{p\leq x} d(p-1)$, proving that 
 $$\sum_{p\leq x} \frac{1}{d(p-1)} \asymp \frac{x}{(\log x)^{3/2}}.$$ This result, together with Erd\H os's asymptotic formula 
 $ \sum_{n \leq x} d(d(n)) \sim c \, x \log \log x $ for a constant $c \in (0, \infty)$, lead us to consider the hybrid $\sum_{n \leq x} 
 \frac{1}{d(d(n))}$ of the Erd\H os and Ramanujan divisor sums. The presence of the reciprocal significantly complicates the analysis, as it 
 amplifies the contribution of integers for which $d(d(n))$ is exceptionally small. In this paper, we prove that $$\sum_{n \leq x} 
 \frac{1}{d(d(n))} \asymp \frac{x}{ \log \log x}, $$ through a combined application of Golomb's estimate for powerful numbers 
 and Tur\'{a}n's quantitative form of the Hardy--Ramanujan theorem. 
\end{abstract}

\vspace{0.1in}

\noindent {\footnotesize \emph{MSC:} 11N37}

\vspace{0.1in}

\noindent {\footnotesize \emph{Keywords:} divisor, divisor sum, prime divisor, prime omega function, Titchmarsh divisor problem, 
 squarefree part}

\section{Introduction}
 The divisor function $d(n) = \sum_{d \mid n} 1 $ is among the most fundamental arithmetic functions. 
 The iteration 
 $d(d(n))$ 
 is closely related to Ramanujan's seminal work on highly composite numbers \cite{Ramanujan1915comp}, and this motivated 
 Erd\H{o}s~\cite{Erdos1968} to investigate the behavior of sums of the form $ \sum_{n \leq x} d(d(n))$. 
 In this direction, Erd\H os proved 
 the remarkable result whereby: There exists a constant $c_1$ such that $ 0 < c_1 < \infty$ and such that 
\begin{equation}\label{displayErdos}
 \sum_{n \leq x} d(d(n)) \sim c_1 \, x \log \log x. 
\end{equation}
 This may be seen in contrast to Dirichlet's formula 
\begin{equation}\label{displayDirichlet}
 \sum_{n\leq x} d(n) = x \log x + (2 \gamma -1) x + O(\sqrt{x}), 
\end{equation}
 for the Euler--Mascheroni constant $\gamma = \lim_{n \to \infty} \big( 1 + \frac{1}{2} 
 + \cdots + \frac{1}{n} - \log n \big)$, in the sense that Dirichlet’s hyperbola method 
 used in the classical derivation of \eqref{displayDirichlet} cannot be applied directly to the sum in \eqref{displayErdos}. 
 Since $d(d(n))$ is not multiplicative, in contrast to $d(n)$ being multiplicative, this may be seen 
 as complicating the analysis of the behavior of $d(d(n))$. 
 Indeed, the difficulties surrounding the behavior of the self-composition $d(d(n))$ may be seen by 
 the solution due to Buttkewitz et al.\ in 2012 \cite{ButtkewitzElsholtzFordSchlagePuchta2012} of a problem due to Ramanujan 
 \cite{Ramanujan1915comp} concerning the maximal order, asymptotically, of $\log d(d(n))$. The Erd\H os estimate in 
 \eqref{displayErdos} may also be seen in relation to a number of subsequent research works improving upon or otherwise extending 
 this estimate~\cite{ErdosKatai1969,Heppner1974,Katai1969,Rieger1972}. 

 In 1915, Ramanujan \cite{Ramanujan1915analytic} introduced asymptotic evaluations for a number of divisor sums, including 
\begin{multline}\label{displayRamanujan}
 \sum_{n \leq x} \frac{1}{d(n)} = \\ 
 x \left( \frac{A_1}{ (\log x)^{\frac{1}{2}} } 
 + \frac{A_2}{ (\log x)^{\frac{3}{2}} } 
 + \cdots + \frac{A_r}{(\log x)^{r - \frac{1}{2}}} + O\left( \frac{1}{(\log x)^{r + \frac{1}{2}}} \right) \right), 
\end{multline}
 for $$ A_{1} = \frac{1}{\sqrt{\pi}} \prod_p \sqrt{p^2 - p} \log\left( \frac{p}{p-1} \right) $$ and for unspecified constants $A_{2}$, $A_{3}$, 
 $\ldots$, $A_{r}$, with a full proof of Ramanujan's estimate having been introduced by Wilson in 1923 \cite{Wilson1923}. If we compare 
 the Erd\H os divisor sum $ \sum_{n \leq x} d(d(n))$ to the Ramanujan divisor sum $ \sum_{n \leq x} \frac{1}{d(n)}$, this 
 and the work of Gabdullin et al.\ \cite{GabdullinKonyaginIudelevich2023} reviewed below 
 raise questions concerning the 
 hybrid divisor sum $\sum_{n \leq x} \frac{1}{d(d(n))}$, which provides the main object of study this paper. 

\subsection{A hybrid Titchmarsh--Ramanujan divisor sum}
 The \emph{Titchmarsh divisor problem} or \emph{Titchmarsh--Linnik divisor problem} \cite{Halberstam1967} refers to the problem given 
 by determining the asymptotic behavior of the divisor sum $\sum_{p \leq x} d(p-1)$, with the estimate 
\begin{equation}\label{firstTitchmarsh}
 \sum_{p \leq x} d(p-1) = O(x) 
\end{equation}
 having been introduced and proved by Titchmarsh in 1930 \cite{Titchmarsh1930}. Letting $\zeta(x)$ denote the Riemann zeta function, 
 a remarkable improvement to \eqref{firstTitchmarsh} was introduced and proved in 1961 by Linnik \cite{Linnik1961}, who proved, via a 
 dispersion method, that 
\begin{equation}\label{displayLinnik}
 \sum_{p \leq x} d(p-1) \sim \frac{ \zeta(2) \zeta(3) }{ \zeta(6) } x 
\end{equation}
 with an error term $\ll \frac{x}{ (\log x)^{\alpha} }$ for $0 < \alpha < 1$, and one may compare this against simplified proofs of this 
 estimate due to Rodriquez \cite{Rodriquez1965} and Halberstam \cite{Halberstam1967}, along with the improvement upon Linnik's 
 error term due to Bombieri et al.~\cite{BombieriFriedlanderIwaniec1986}. 

 As a natural hybrid of the Ramanujan divisor sum $ \sum_{n \leq x} \frac{1}{d(n)}$ and the Titchmarsh divisor sum shown in 
 \eqref{firstTitchmarsh} and \eqref{displayLinnik}, 
 Gabdullin et al.\ \cite{GabdullinKonyaginIudelevich2023}
 studied the divisor sum 
 $\sum_{p\leq x} \frac{1}{d(p-1)}$ and succeeded in proving 
 the order of magnitude evaluation 
\begin{equation}\label{magnitudeGabdullin}
 \sum_{p\leq x} \frac{1}{d(p-1)} \asymp \frac{x}{(\log x)^{3/2}}, 
\end{equation}
 thus improving upon a previous result due to Iudelevich \cite{Iudelevich2022}. Gabdullin et al.\ also conjectured that there exists a 
 constant $c_2$ such that $0 < c_2 < \infty$ and such that 
\begin{equation}\label{Gabdullinconj}
 \sum_{p\leq x} \frac{1}{d(p-1)} \sim c_{2} \frac{x}{(\log x)^{3/2}}
\end{equation}
 and noted that this is likely a difficult conjecture. 

\subsection{A hybrid Erd\H{o}s--Ramanujan divisor sum}
 The problem of estimating the divisor sum indicated in \eqref{magnitudeGabdullin} and \eqref{Gabdullinconj} is referred to as 
 \emph{Karatsuba’s divisor problem} \cite{GabdullinKonyaginIudelevich2023,Iudelevich2022}. Since this concerns the hybrid $\sum_{p 
 \leq x} \frac{1}{d(p-1)}$ of the Ramanujan divisor sum $\sum_{n\leq x} \frac{1}{d(n)}$ and the Titchmarsh divisor sum $\sum_{p \leq 
 x} d(p - 1)$, this motivates the study of the corresponding hybrid divisor sum 
 $\sum_{n \leq x} \frac{1}{d(d(n))}$ given by 
 a ``fusion'' of the same Ramanujan divisor sum 
 and the Erd\H os divisor sum $\sum_{n \leq x} d(d(n))$. 

 We succeed in proving that $$\sum_{n \leq x} \frac{1}{d(d(n))} \asymp \frac{x}{ \log \log x}, $$ by analogy with both the 
 Gabdullin--Konyagin--Iudelevich order of magnitude estimate in \eqref{magnitudeGabdullin} and the Erd\H os asymptotic estimate in 
 \eqref{displayErdos}. We also conjecture, by analogy with the   conjecture in \eqref{Gabdullinconj}, that there exists a constant  $c_{3}$ such  
  that $0 < c_{3} < \infty$ and such that 
\begin{equation}\label{newconj}
 \sum_{n \leq x} \frac{1}{d(d(n))} \sim c_{3} \frac{x}{ \log \log x}. 
\end{equation}
 As is the case with the Gabdullin--Konyagin--Iudelevich conjecture in   \eqref{Gabdullinconj}, the conjecture in \eqref{newconj} seems to 
  be difficult. 
 
 Informally, when dealing with sums of reciprocals of divisors, as opposed to ``non-reciprocal'' divisor sums, applying the reciprocal 
 operation on divisors can be thought of as accentuating  the irregular behavior of $d(n)$ or its subsequences.  This is evidenced by 
 Gabdullin et al.\ \cite{GabdullinKonyaginIudelevich2023} commenting on the difficulty of the conjecture in  \eqref{Gabdullinconj}, in  
 contrast to the original  Titchmarsh divisor problem. 
 This is also reflected by the contrast between the estimate for Dirichlet's original divisor sum in 
 \eqref{displayDirichlet} compared with the estimate
 for Ramanujan's reciprocal divisor sum in \eqref{displayRamanujan}. 

 The research contribution due to Ramanujan \cite{Ramanujan1915analytic} that introduced the divisor sum relation in 
 \eqref{displayRamanujan} may be seen as seminal in the history of divisor sums. 
 In this same contribution, Ramanujan introduced
 an asymptotic evaluation for 
 $ \sum_{n \leq x} d^{2}(n)$, 
 which has inspired a number of notable 
 advances on divisor sums of this form \cite{CullyHugillTrudgian2021,JiaSankaranarayanan2014,MaierSankaranarayanan2005,RamachandraSankaranarayanan2003}. 

\section{Required preliminaries}\label{secprelim}
 For real-valued functions $f$ and $g$ (defined, say, on a subset of $\mathbb{R}_{\geq 0}$), if there exists a constant $C > 0$ and a value 
 $x_0$ such that $ |f(x)| \leq C g(x) $ for all $x \geq x_0$, then we write $f(x) \ll g(x)$. Observe that this is equivalent to $f(x) = O(g(x))$. 
 In a similar spirit, if both $f(x) \ll g(x)$ and $g(x) \ll f(x)$, then it is standard to write $f(x) \asymp g(x)$. 

 The prime omega function $\omega(n)$ gives the number of distinct prime factors of $n \geq 2$, with $\omega(1) = 0$. Similarly, we 
 have that $\Omega(n)$
 gives the total number of prime factors of $n$, counting multiplicities. 
 We then let \(a(n)\) denote the product of the primes \(p\) such that \(p\Vert n\), 
 i.e., the \emph{squarefree} part of \(n\) supported on primes dividing \(n\) exactly once. 
 By letting the prime factorization of $n$ be written as
 $n = \prod_{i} p_{i}^{e_i}$, the \emph{powerful part} of $n$
 may be defined so that 
\begin{equation}\label{displaybn}
 b(n) = \prod_{e_{i} >1} p_{i}^{e_{i}}. 
\end{equation} 
 The definitions of $d(n)$, $\omega(n)$, $a(n)$, and $b(n)$ together give us that 
\begin{equation}\label{keyErdos}
 d(n) = 2^{\omega(a(n))} d(b(n)), 
\end{equation}
 and \eqref{keyErdos} provides a key in the approach used by Erd\H os~\cite{Erdos1968} in the derivation of the 
 asymptotic relation on display in \eqref{displayErdos}. 

\subsection{The Abel summation formula}
 Since this paper concerns average values of arithmetic functions, it is natural that summation techniques often employed in analytic 
 number theory would arise. For a full preliminary treatment of such summation tools, we refer to the appropriate section of De Koninck 
 and Luca's text on analytic number theory \cite[\S1]{DeKoninckLuca2012}, and the following formulation of the Abel summation lemma is 
 required for our purposes. 
 
 \ 

\noindent {\bf (Abel summation formula):} Let $( a_{n} )_{n \geq 1}$ be a sequence of complex numbers and let $f : [1, +\infty) \to 
 \mathbb{C}$ have continuous derivatives for $x \geq 1$. For each real number $x \geq 1$, let $$ A(x) = \sum_{n \leq x} a_n. $$ Then 
 \begin{equation}\label{eq1p4} 
 \sum_{n \leq x} a_{n} f(n) = A(x) f(x) - \int_{1}^{x} A(t) f'(t) \, dt. 
 \end{equation}

\subsection{Golomb's bounds for powerful numbers}
 A positive integer $n$ is said to be \emph{powerful} if $p \mid n \Longrightarrow p^2 \mid n$ for an arbitrary prime $p$. Equivalently, 
 we have that $n = b(m)$ for some $m$, according to the definition in \eqref{displaybn}. Being consistent with notation from Golomb 
 \cite{Golomb1970}, we write $k(x)$ in place of the number of powerful numbers $\leq x$. A remarkable result due to Golomb is such that 
\begin{equation}\label{displayGolomb}
 c_4 x^{1/2} - 3 x^{1/3} \leq k(x) \leq c_4 x^{1/2} 
\end{equation}
 for the constant $c_4 = \frac{\zeta\left( \frac{3}{2} \right)}{\zeta(3)}$. We apply this result from Golomb in Section \ref{secmain} below. 

\subsection{Tur\'an's proof of the Hardy--Ramanujan theorem}\label{secTuran}
 The \emph{Tur\'an--Kubilius inequality} provides a powerful tool for studying the average behavior of additive arithmetic functions 
 \cite[\S3.2]{Tenenbaum2015}. Being consistent with notation from Tenenbaum's text on analytic and probabilistic number theory 
 \cite[\S3.2]{Tenenbaum2015}, for an arithmetic function $f$, we denote the associated mean/expectation so that 
\begin{equation}\label{eq3p2}
 g(N) = \mathbb{E}_{N}(f) := \frac{1}{N} \sum_{1 \leq n \leq N} f(n). 
\end{equation} 
 As suggested by Tenenbaum, one would hope, depending on the input function $f$, that the mean in \eqref{eq3p2} would provide a 
 reasonable approximation to a normal order of $f$. Letting the associated variance be written as 
\begin{equation}\label{eq3p4}
 \mathbb{V}_{N}(f) := \mathbb{E}_{N}\big( |f(n) - g(N)|^2 \big), 
\end{equation} 
 one would want to obtain a relation involving \eqref{eq3p2} and \eqref{eq3p4} to guarantee or formalize how $g$ is, as desired, a 
 normal order for $f$. 
 This provides a way, for the $f(n) = \omega(n)$ case, 
 of proving that 
\begin{equation}\label{displayTuran}
 \sum_{n \leq N} \big( \omega(n) - \log \log N \big)^{2} \ll N \log \log N, 
\end{equation}
 with \eqref{displayTuran} having been proved in a non-probabilistic and elementary way by Tur\'an in 1934 \cite{Turan1934} to prove the 
 Hardy--Ramanujan theorem, which concerns the normal order of $\omega(n)$. Tur\'an's result in \eqref{displayTuran} is to play a 
 central role in our derivation of our main result in Section \ref{secmain} below. 
 
\section{Main result}\label{secmain}

\begin{lemma}\label{lemmaB1}
 Letting $$ \mathcal{B}_{1}(x) := \left\{ n \leq x : b(n) > (\log \log x)^{2} \right\}, $$ we have that $$ \sum_{n \in \mathcal{B}_{1}(x)} 
 \frac{1}{d(d(n))} \ll \frac{x}{\log \log x}. $$
\end{lemma}

\begin{proof} 
 Write $Y = (\log\log x)^2$. For a fixed powerful number $\beta$, the number of positive integers $n \leq x$ such that $b(n) = \beta$ is 
 bounded above by the \emph{total} number $\lfloor \frac{x}{\beta} \rfloor$ of positive multiples of $\beta$ not exceeding $x$, and 
 hence the relations whereby 
\begin{equation}\label{firstA2proof}
\#\mathcal{B}_1(x) \leq \sum_{\substack{\beta > Y\\ \beta \ {\rm powerful}}} \#\{n\le x : b(n) = \beta\} \leq 
 x \sum_{\substack{\beta > Y\\ \beta \ {\rm powerful}}}\frac1\beta.
\end{equation}
 By the Golomb estimate in \eqref{displayGolomb}, we find that 
\begin{equation}\label{consequenceGolomb}
k(t)\ll t^{1/2}.
\end{equation}
 In the formulation of the Abel summation formula given in Section \ref{secprelim}, we set $a_n = 1$ if $n$ is powerful and $a_n = 
 0$ otherwise, and we set $f(t) = \frac{1}{t}$, so that $$ \sum_{\substack{Y<b\le Z\\ b\ {\rm powerful}}} \frac{1}{b} = \frac{k(Z)}{Z} - 
 \frac{k(Y)}{Y} + \int_Y^Z \frac{k(t)}{t^2} \, dt, $$ and hence 
\begin{equation}\label{withoutY}
 \sum_{\substack{Y < b \leq Z\\ b\ {\rm powerful}}}\frac{1}{b} \leq \frac{k(Z)}{Z} + \int_Y^Z \frac{k(t)}{t^2}\,dt. 
\end{equation}
 Using the consequence in \eqref{consequenceGolomb} of the Golomb bounds, and letting \(Z\to\infty\) in \eqref{withoutY} and applying 
 the monotone convergence theorem, it follows that
\begin{equation}\label{aftermonotone}
 \sum_{\substack{b>Y\\ b\ {\rm powerful}}}\frac{1}{b} \ll Y^{-1/2}.
\end{equation} 
 From \eqref{firstA2proof} and \eqref{aftermonotone} together, we find that $ \#\mathcal{B}_1(x)\ll \frac{x}{\log\log x}$. Finally, since 
 $d(d(n)) \geq 1$, we have $$ \sum_{n \in \mathcal{B}_1(x)} \frac1{d(d(n))} \leq \#\mathcal{B}_1(x) \ll \frac{x}{\log\log x}, $$ as desired. 
\end{proof}

\begin{lemma}\label{LemmaB2}
Let $$ \mathcal{B}_2(x) := \left\{n \leq x:\omega(n) \notin \left[\frac12\log\log x,2\log\log x\right]\right\}. $$ Then $$ \sum_{n \in 
 \mathcal{B}_2(x)} \frac{1}{d(d(n))} \ll \frac{x}{\log\log x}. $$
\end{lemma}

\begin{proof}
 If \(n\in\mathcal{B}_2(x)\), then
\begin{equation}\label{firstproofB2}
 \frac12\log\log x \leq \left|\omega(n)-\log\log x\right|. 
\end{equation}
 From  \eqref{firstproofB2}, we   obtain that
\begin{align*}
 \frac{1}{4} (\log \log x)^{2} \# \mathcal{B}_{2}(x) 
 & \leq \sum_{n \in \mathcal{B}_{2}(x)} \big( \omega(n) - \log \log x \big)^{2} \\ 
 & \leq \sum_{n \leq x} \big( \omega(n) - \log \log x \big)^{2} \\ 
 & \ll x \log \log x, 
\end{align*}
 where we have applied the Tur\'{a}n estimate reviewed in Section \ref{secTuran}. So, since \(0 < \frac{1}{d(d(n))} \leq 1\), we see that 
$$ \sum_{n\in\mathcal{B}_2(x)}
\frac1{d(d(n))}
\leq \#\mathcal{B}_2(x)
\ll \frac{x}{\log\log x}, $$
 as desired. 
\end{proof}

\begin{lemma}\label{LemmaComplement}
We have that 
$$ \sum_{\substack{n\leq x\\
n\notin \mathcal{B}_1(x) \cup \mathcal{B}_2(x)}}
\frac1{d(d(n))} \ll
\frac{x}{\log\log x}. $$
\end{lemma}

\begin{proof}
 Let \(n \le x\), with \(n \notin \mathcal{B}_1(x) \cup \mathcal{B}_2(x)\). Recalling the divisor function identity in \eqref{keyErdos}, and 
 following a similar approach as in the work of Erd\H{o}s \cite{Erdos1968}, we write $ d(b(n)) = 2^{u(n)} v(n)$ for $v(n)$ odd, giving 
 us that 
\begin{equation}\label{userelative}
 d(n) = 2^{\omega(a(n)) + u(n)} v(n).
\end{equation}
 Since $2^{\omega(a(n)) + u(n)}$ and $v(n)$ are relatively prime, we obtain from \eqref{userelative} that $ d(d(n)) = \big(\omega(a(n)) + 
 u(n) + 1\big) d(v(n))$. This gives us the lower bound 
\begin{equation}\label{lowerboundddn}
d(d(n)) \ge \omega(a(n)) + 1.
\end{equation}
 We proceed to estimate \(\omega(a(n))\). Since \(n\notin \mathcal B_2(x)\), we have that 
\begin{equation}\label{consequencenotB2}
\omega(n)\geq \frac12\log\log x.
\end{equation}
 Moreover, since \(n\notin \mathcal B_1(x)\), we have that 
\begin{equation}\label{sincenotinB1}
 b(n) \leq (\log\log x)^2.
\end{equation}
 Since 
\begin{equation}\label{omegadiff}
 \omega(a(n))=\omega(n)-\omega(b(n)), 
\end{equation}
 and since \eqref{sincenotinB1} gives us that 
$ 2^{\omega(b(n))} \leq b(n) \leq (\log\log x)^2$,
we obtain
\begin{equation}\label{omegalogloglog}
 \omega(b(n))\leq \frac{2\log\log\log x}{\log 2}.
\end{equation}
 A combined application of \eqref{consequencenotB2}, \eqref{omegadiff}, and \eqref{omegalogloglog} then allows us to obtain that 
\begin{equation}\label{omegaangeq}
 \omega(a(n)) \geq \frac12\log\log x - \frac{2\log\log\log x}{\log 2} \gg \log\log x.
\end{equation}
 In turn, combining \eqref{lowerboundddn} and \eqref{omegaangeq} allows us to deduce that 
\begin{equation}\label{notAbel}
 \frac{1}{d(d(n))} \ll \frac{1}{\log\log x}, 
\end{equation}
 uniformly for $n \leq x$ such that $n \not\in \mathcal{B}_{1}(x) \cup \mathcal{B}_{2}(x)$. Applying the summation operator 
 $\sum_{\substack{n \leq x \\ n \not\in \mathcal{B}_{1}(x) \cup \mathcal{B}_{2}(x) }} \cdot $ to both sides of \eqref{notAbel} then allows 
 us to obtain the desired result. 
\end{proof}

\begin{lemma}\label{LemmaLower}
 We have that $$ \frac{x}{\log\log x} \ll \sum_{n\le x}\frac{1}{d(d(n))}. $$
\end{lemma}

\begin{proof}
Let
\begin{equation}\label{displayS2}
\mathcal{B}_{3}(x) := \left\{ n \le x : n \text{ squarefree},\ 
\omega(n) \in \left[\frac12\log\log x, 2\log\log x\right] \right\}.
\end{equation}
 It is classically and well known that $ \#\{n \le x : n \text{ squarefree}\} = \frac{6}{\pi^2}x + O(\sqrt{x})$. 
 By rewriting \eqref{displayS2} as a set difference involving $\mathcal{B}_{2}$, we find that
\begin{equation}\label{fromsetdiff}
 \#\mathcal{B}_{3}(x) \geq \#\{ n \leq x : \text{$n$ squarefree} \} 
 - \#\mathcal{B}_{2}(x). 
\end{equation}
 It was shown in the proof of Lemma \ref{LemmaB2} that $ \#\mathcal{B}_{2}(x) \ll \frac{x}{\log \log x}$. This together with the classical 
 estimate for $\#\{n \le x : n \text{ squarefree}\}$ then give us, from \eqref{fromsetdiff}, that $\#\mathcal{B}_{3}(x) \gg x$. Now, let 
 $n \in \mathcal{B}_{3}(x)$. Since \(n\) is squarefree, we have that the relation $ d(n) = 2^{\omega(n)}$ holds. As a consequence, we 
 find that the equalities such that $ d(d(n)) = d\big(2^{\omega(n)}\big)=\omega(n)+1 $ hold true. Moreover, the definition of 
 $\mathcal{B}_{3}(x)$ gives us that $ \omega(n)+1\leq 2\log\log x+1 \ll \log\log x$, so that $ \frac{1}{d(d(n))} = \frac{1}{\omega(n) + 
 1} \gg \frac{1}{\log\log x}$. It then follows that $$ \sum_{n \leq x} \frac{1}{d(d(n))} \geq \sum_{n\in\mathcal{B}_{3}(x)}\frac{1}{d(d(n))}
\gg \frac{\#\mathcal{B}_{3}(x)}{\log\log x} \gg
\frac{x}{\log\log x}, $$ giving us the desired result. 
\end{proof}

\begin{theorem}\label{MainTheorem}
We have
\begin{equation}\label{displaymain}
\sum_{n\le x}\frac{1}{d(d(n))}
\asymp
\frac{x}{\log\log x}.
\end{equation}
\end{theorem}

\begin{proof}
 The upper bound associated with the order of magnitude estimate in \eqref{displaymain} follows by decomposing the index set 
 $$ \{ n : n \leq x \} = \mathcal{B}_1(x) \cup \mathcal{B}_2(x) \cup \big( \{ n : n \leq x \} \setminus(\mathcal{B}_1(x) \cup 
 \mathcal{B}_2(x))\big) $$ and applying Lemmas \ref{lemmaB1}, \ref{LemmaB2}, 
 and \ref{LemmaComplement}. The corresponding lower bound is given in Lemma \ref{LemmaLower}.
\end{proof}

\section{Conclusion}
 Recall Gabdullin et al.\ \cite{GabdullinKonyaginIudelevich2023} commented that the conjectured asymptotic equivalence in 
 \eqref{Gabdullinconj} appears to be difficult. It appears that the conjectured asymptotic equivalence corresponding to our main result 
 in Theorem \ref{MainTheorem} is similarly difficult, and we conclude by encouraging the exploration of this. 

\subsection*{Acknowledgements}
 The author is very thankful to an anonymous referee for a previous, unpublished version of this work for a number of very helpful 
 suggestions. Notably, the decomposition of integers into squarefree and powerful parts and the treatment of exceptional sets used in 
 this paper are partly based on the referee’s report. The author is also very thankful to Karl Dilcher for many useful discussions related to 
 this paper and thanks Marco Cantarini and Amiram Eldar for useful commentary related to this paper. 
 The author used the GPT-5.3 model to be of assistance with exposition and to explore heuristic arguments, but 
 all mathematical results and proofs were verified independently by the author.

\subsection*{Statements and Declarations}
 There are no financial or non-financial competing interests that are directly or indirectly related to the work 
 to disclose. 

 \bibliographystyle{plain}
\bibliography{seprefe}

\end{document}
